\section{Steps with a common singular series}\label{sec:sparse}

Write \(S_w(N)=L^2X_N/N\) for the truncated singular series in
\eqref{eq:parameters}. We can select integer multipliers for which a
growing initial segment of the progression has this same factor to
arbitrary inverse-logarithmic accuracy.

\begin{theorem}\label{thm:sparse}
Fix \(A,c>0\) and set \(r=\lfloor cw/\log w\rfloor\).
For every sufficiently large even \(N\) and every admissible \(H\),
at least \(Y L^{-c-o(1)}\) integers \(k\in(Y,2Y]\) satisfy
\begin{equation}\label{eq:sparse}
    G(n_i)=n_i S_w(N)+O_{A,c,\varepsilon}(N/L^A)
    \qquad (1\le i\le r).
\end{equation}
\end{theorem}

\begin{proof}
Put \(z=L^{A+4}\). For each prime \(w<p\le z\), exclude the residues
of \(k\bmod p\) for which \(p\mid N-iWk\) for at least one
\(1\le i\le r\). For large \(N\) we have \(r<w/2\). Thus \(iW\) is
invertible modulo \(p\), and the number of excluded residues is
\[
 \rho(p)=
 \begin{cases}
  1,&p\mid N,\\
  r,&p\nmid N .
 \end{cases}
\]
Indeed, the residues \(N(iW)^{-1}\bmod p\) coincide at zero in the
first case and are distinct in the second. Let
\[
 V=\prod_{w<p\le z}(1-\rho(p)/p),\qquad
 U=\sum_{w<p\le z}\rho(p)/p .
\]
By Mertens' estimate,
\(U\le r(\log\log z-\log\log w+O(1))\le(c+o(1))w\).
Also \(\sum_{w<p\le z}(\rho(p)/p)^2\ll r^2/w=o(w)\). Expanding the
logarithm gives \(\log V\ge-(c+o(1))w\).

We next count the surviving integers in \((Y,2Y]\). Fix \(B>c+1\).
For a squarefree product \(d\) of the sieving primes, put
\(\rho(d)=\prod_{p\mid d}\rho(p)\). The Chinese remainder theorem gives
\(Y\rho(d)/d+O(\rho(d))\) multipliers in the intersection of the
corresponding excluded classes. Truncate inclusion--exclusion at the
consecutive orders \(t-1,t\), with \(t\) odd. The two Bonferroni bounds
for the model proportion bracket \(V\) and differ by at most
\(U^t/t!\). Choose \(t=Cw+O(1)\) with \(C\) sufficiently large in
terms of \(A,c,B\). Then \(U^t/t!\le(eU/t)^t\ll L^{-B}\).
For either truncation the total rounding error, divided by \(Y\), is
at most a constant times
\[
 Y^{-1}\!\!\sum_{\substack{d\\\omega(d)\le t}}\rho(d)
 \le Y^{-1}(1+rz)^t
 =Y^{-1}\exp(O_{A,c,B}(w^2))=o(L^{-B}),
\]
where the sum runs over the squarefree products just specified and
\(\log Y\gg_\varepsilon L\). The integer count is therefore
\(Y(V+O_{A,c,B,\varepsilon}(L^{-B}))\).

For a surviving \(k\), every prime divisor of \(n_i\) up to \(z\)
already lies below or at \(w\), where \(n_i\equiv N\pmod W\) fixes
the local factors. For all \(i\le r\) we have
\[
 0\le\log\frac{\SSer(n_i)}{S_w(N)}
   \ll\sum_{\substack{p>z\\p\mid n_i}}\frac1p
   \le \frac{L}{z\log z}.
\]
Mertens' product estimate gives \(S_w(N)\ll\log w\).
Thus \(n_i\SSer(n_i)=n_i S_w(N)+O_A(N/L^A)\).

Finally apply \eqref{eq:mrt} on \([N-2iH,N-iH]\) for each \(i\le r\).
As in Lemma~\ref{lem:transfer}, the exceptional multipliers number
\(O_{E,\varepsilon}(r^2H/L^E)\). Since \(W=L^{1+o(1)}\) and \(H=WY\),
we may choose \(E>A+B+3\) to bound this number by \(O(YL^{-B})\).
Removing higher prime powers costs \(O(\sqrt N L^3)\) at each target.
The resulting count of multipliers is still at least \(Y L^{-c-o(1)}\),
and each satisfies \eqref{eq:sparse}.
\end{proof}

When \(N\) runs through powers of two, each sieving prime excludes exactly
\(r\) classes. In this case \(U=(c+o(1))w\) and the sieve itself retains
\(Y L^{-c+o(1)}\) multipliers. This measures the size of the particular
selection used in the proof.
